\documentclass[12pt,a4paper]{report}
\usepackage[dvips]{graphicx}           
\usepackage[portuges]{babel}           
\usepackage{color}            
\usepackage[latin1]{inputenc}           
\usepackage[normalem]{ulem}
\setlength{\textwidth}{160mm}
\setlength{\textheight}{220mm}
\setlength{\topmargin}{-8pt}
\setlength{\oddsidemargin}{0mm}
\usepackage{amsmath,amsfonts, amssymb}
\usepackage{enumitem}% http://ctan.org/pkg/enumitem

\begin{document}

  \begin{minipage}[b]{0.1 \linewidth}
     \includegraphics[width=.9cm]{ufs_vertical_positiva.png}
  \end{minipage}\hfill	
  \begin{minipage}[b]{0.85 \linewidth}
	       \textsc{\large Universidade Federal de Sergipe}\\
	       \textsc{\large Departamento de Estat�stica e Ci�ncias Atuariais}\\
	       Disciplina: ESTAT0078 -- Infer�ncia I\\
	       Professor: Sadraque E. F. Lucena
\end{minipage}

\vspace{0.8cm}
\begin{center} {\bf Lista de Exerc�cios 11} \end{center}
\vspace{0.8cm}

	\begin{enumerate}[label=11.\arabic*)]
		
		% Quest�o 1
		\item Seja $X_1,\ldots,X_n$ uma amostra aleat�ria da vari�vel aleat�ria $X$. Encontre  estat�sticas conjuntamente suficientes considerando as distribui��es a seguir:
		\begin{enumerate}
			\item $X\sim$ Gama Inversa($\alpha,\beta$): $f(x) = \frac{\beta^\alpha}{\Gamma(\alpha)} x^{-(\alpha+1)} e^{-\beta/x}$, $x>0$, $\alpha,\beta>0$
			\item $X\sim$ Beta($\alpha,\beta$): $f(x) = \frac{\Gamma(\alpha+\beta)}{\Gamma(\alpha)\Gamma(\beta)} x^{\alpha-1}(1-x)^{\beta-1}$, $0<x<1$, $\alpha,\beta>0$
		\end{enumerate}
		
	\end{enumerate}
	
	\newpage
	
	\noindent Respostas:
	
	\begin{enumerate}[label=11.\arabic*)]
		
		% Quest�o 1
		\item 
		\begin{enumerate}
			\item $\left(\prod\limits_{i=1}^n x_i, \sum\limits_{i=1}^n x_i \right)$ s�o conjuntamente suficientes para $(\alpha,\beta)$.
			
			\item $\left(\prod\limits_{i=1}^n x_i, \prod\limits_{i=1}^n (1-x_i) \right)$ s�o conjuntamente suficientes para $(\alpha,\beta)$.
		\end{enumerate}
		
	\end{enumerate}
	
\end{document}